% Univalence + HITs (Ch 3-4)
% Lines 560-763 of main.tex
% This file is for agent editing — will be merged back

\chapter{The Univalence Axiom}
\label{ch:univalence}

\section{Type Equivalences}

\begin{definition}[Equivalence of Types]
A function $f : A \to B$ is an \emph{equivalence} if there exists
$g : B \to A$ with:
\begin{align}
  &\prod_{a : A} \Path_A(g(f(a)), a) \\
  &\prod_{b : B} \Path_B(f(g(b)), b)
\end{align}
We write $A \simeq B$ if there exists such an $f$.
\end{definition}

\section{The Univalence Axiom}

\begin{axiom}[Univalence]
\label{ax:univalence}
For types $A, B : \Type$, the canonical map
\[
  \mathsf{idToEquiv} : (A =_\Type B) \to (A \simeq B)
\]
is itself an equivalence.
\end{axiom}

In cubical type theory, univalence is a \emph{theorem}, not an axiom:
it is derived from the Kan structure and the glueing construction.

\section{Univalence for Program Semantics}

For our purposes, univalence has two computational consequences:

\begin{enumerate}
  \item \textbf{Shared function symbols}: If two programs both call
    Python's \texttt{sorted}, these references denote the \emph{same}
    function.  In the universe $\Type$, the type ``Python's sorted
    function'' has a unique representative.  By univalence, any two
    copies of this type are equal.  Computationally, we implement this
    as a shared Z3 function symbol: both programs use
    $\mathsf{py\_sorted}$.

  \item \textbf{Type coercions as paths}: Python's type coercions
    (e.g., $\mathsf{bool} \to \mathsf{int}$ via \texttt{int(True)} = 1)
    are \emph{equivalences} between types.  By univalence, they become
    \emph{paths} in $\Type$.  Transport along these paths relates
    the behaviors of programs that use different but coercible types.
\end{enumerate}

\begin{example}[Shared \texttt{sorted}]
Program $f$ contains \texttt{sorted(x)} and program $g$ contains
\texttt{sorted(y)}.  In the Z3 model, both compile to
$\mathsf{py\_sorted}(\den{x})$ and $\mathsf{py\_sorted}(\den{y})$
respectively --- the \emph{same} Z3 function.  Z3 can then focus on
whether $\den{x} = \den{y}$, rather than comparing two independent
uninterpreted functions.
\end{example}

\begin{example}[Bool--Int Coercion Path]
The coercion $c : \mathsf{bool} \simeq \{0, 1\} \hookrightarrow
\mathsf{int}$ gives a path $p : \Path_\Type(\mathsf{bool},
\mathsf{int})$ (in the subuniverse of types with $\leq 2$ values).
If program $f$ returns a \texttt{bool} and program $g$ returns
\texttt{0} or \texttt{1}, transport along $p$ relates them.
\end{example}


\chapter{Higher Inductive Types for Programs}
\label{ch:hits}

\section{Definition and Motivation}

A \emph{higher inductive type} (HIT) is a type defined by:
\begin{enumerate}
  \item \textbf{Point constructors}: elements of the type.
  \item \textbf{Path constructors}: equalities between elements.
  \item \textbf{Higher path constructors}: equalities between
    equalities (optional --- for coherence).
\end{enumerate}

For program equivalence, we define the HIT of Python computations.
The point constructors are the syntactic forms (literals, operations,
conditionals, loops, etc.).  The path constructors are the
\emph{equational axioms} that identify syntactically different but
semantically equivalent expressions.

\section{The HIT $\mathsf{PyComp}$}

\begin{definition}[The HIT of Python Computations]
\label{def:pycomp}
$\mathsf{PyComp}$ is defined by:

\medskip
\textbf{Point constructors:}
\begin{align}
  \mathsf{lit} &: \Int \to \mathsf{PyComp} \\
  \mathsf{binop} &: \mathsf{OpTag} \times \mathsf{PyComp} \times
    \mathsf{PyComp} \to \mathsf{PyComp} \\
  \mathsf{unop} &: \mathsf{OpTag} \times \mathsf{PyComp} \to
    \mathsf{PyComp} \\
  \mathsf{cond} &: \mathsf{PyComp} \times \mathsf{PyComp} \times
    \mathsf{PyComp} \to \mathsf{PyComp} \\
  \fold &: \mathsf{OpTag} \times \mathsf{PyComp} \times \Nat \times
    \Nat \to \mathsf{PyComp} \\
  \mathsf{call} &: \mathsf{FnSym} \times \mathsf{PyComp}^* \to
    \mathsf{PyComp} \\
  \mathsf{recfn} &: \mathsf{FnSym} \times \mathsf{PyComp}^* \to
    \mathsf{PyComp}
\end{align}

\textbf{Path constructors (the 24 dichotomy axioms, plus arithmetic):}

\emph{Arithmetic paths:}
\begin{align}
  \mathsf{add\_comm} &: \prod_{a,b} \Path(\mathsf{binop}(+, a, b),\;
    \mathsf{binop}(+, b, a)) \\
  \mathsf{mul\_comm} &: \prod_{a,b} \Path(\mathsf{binop}(\times, a, b),\;
    \mathsf{binop}(\times, b, a)) \\
  \mathsf{add\_assoc} &: \prod_{a,b,c} \Path(
    \mathsf{binop}(+, \mathsf{binop}(+, a, b), c),\;
    \mathsf{binop}(+, a, \mathsf{binop}(+, b, c))) \\
  \mathsf{mul\_assoc} &: \prod_{a,b,c} \Path(
    \mathsf{binop}(\times, \mathsf{binop}(\times, a, b), c),\;
    \mathsf{binop}(\times, a, \mathsf{binop}(\times, b, c))) \\
  \mathsf{add\_zero} &: \prod_a \Path(\mathsf{binop}(+, a,
    \mathsf{lit}(0)),\; a) \\
  \mathsf{mul\_one} &: \prod_a \Path(\mathsf{binop}(\times, a,
    \mathsf{lit}(1)),\; a) \\
  \mathsf{mul\_zero} &: \prod_a \Path(\mathsf{binop}(\times, a,
    \mathsf{lit}(0)),\; \mathsf{lit}(0)) \\
  \mathsf{neg\_neg} &: \prod_a \Path(\mathsf{unop}(-,
    \mathsf{unop}(-, a)),\; a) \\
  \mathsf{abs\_idem} &: \prod_a \Path(\mathsf{unop}(|\cdot|,
    \mathsf{unop}(|\cdot|, a)),\; \mathsf{unop}(|\cdot|, a)) \\
  \mathsf{min\_comm} &: \prod_{a,b} \Path(\mathsf{binop}(\min, a, b),\;
    \mathsf{binop}(\min, b, a)) \\
  \mathsf{max\_comm} &: \prod_{a,b} \Path(\mathsf{binop}(\max, a, b),\;
    \mathsf{binop}(\max, b, a))
\end{align}

\emph{Dichotomy paths (one per dichotomy, see Part II):}
\begin{align}
  \mathsf{D1\_rec\_iter} &: \prod_{\mathsf{base}, \mathsf{step}}
    \Path(\rec(\mathsf{base}, \mathsf{step}),\;
    \fold(\mathsf{step}, \mathsf{base})) \\
  \mathsf{D2\_beta} &: \prod_{f, a}
    \Path((\lambda x.\, \mathsf{body})(a),\;
    \mathsf{body}[x \mapsto a]) \\
  \mathsf{D3\_while\_for} &: \prod_{\mathsf{init}, \mathsf{step}, n}
    \Path(\mathsf{while\_loop}(\mathsf{init}, \mathsf{step}, n),\;
    \mathsf{for\_loop}(\mathsf{init}, \mathsf{step}, n)) \\
  &\vdots \quad \text{(24 dichotomies total, see Part II)}
\end{align}

\textbf{Truncation:}
$\mathsf{PyComp}$ is 0-truncated (a set): any two parallel paths
are equal.  This reflects that we care about \emph{whether} programs
are equivalent, not about \emph{how many proofs} there are.
\end{definition}

\section{The Elimination Principle}

\begin{theorem}[Elimination for $\mathsf{PyComp}$]
\label{thm:pycomp-elim}
To define a function $h : \mathsf{PyComp} \to X$ (for a set $X$),
it suffices to give:
\begin{enumerate}
  \item A value $h(\mathsf{lit}(n))$ for each literal.
  \item A value $h(\mathsf{binop}(\op, a, b))$ for each binary operation,
    given $h(a)$ and $h(b)$.
  \item Similarly for all other point constructors.
  \item A proof that $h$ \emph{respects} each path constructor:
    $h(p(\mathsf{i0})) = h(p(\mathsf{i1}))$ for each path $p$.
\end{enumerate}
\end{theorem}

\begin{proof}
This is the standard elimination principle for HITs
(Lumsdaine, 2011; Shulman, 2015).
\end{proof}

The elimination principle means: any function defined on Python
computations that respects the equational axioms factors uniquely
through $\mathsf{PyComp}$.  In particular, the Z3 denotation functor
$\den{\cdot}$ is such a function: it maps syntax to Z3 terms and
respects the axioms (because the axioms are added to the Z3 solver).

\section{The Free Algebra Interpretation}

$\mathsf{PyComp}$ is the \emph{free algebra} generated by the Python
operations modulo the equational axioms:
\[
  \mathsf{PyComp} = \mathsf{Free}_{\Sigma/E}(X)
\]
where $\Sigma$ is the algebraic signature (operations) and $E$ is
the equational theory (path constructors), and $X$ is the set of
program variables.

Two programs are equivalent in $\mathsf{PyComp}$ iff they are equal
in the quotient algebra $\Sigma/E$ --- i.e., connected by a finite
sequence of axiom applications.


